\chapter{数学与概率工具}

\begin{quote}\small
\textit{本章摘要。}本章为后文组织数据、求解参数和描述不确定性准备数学工具。向量与矩阵用于表示多维观测和线性计算。投影、正定性与矩阵分解用于分析解的存在性、唯一性和数值敏感性。微积分说明目标函数怎样随参数变化。概率部分建立条件分布、期望和贝叶斯公式，并说明信息量、期望微分与指数族的计算条件。阅读具体模型时，可以返回相应小节核对记号与假设。
\end{quote}

\section{向量、矩阵与基本运算}
设备的一次观测可能同时包含温度、振动和转速。向量把同一次观测中的数值排在一起，矩阵把多次观测或多条计算规则排在一起。本节先说明这些对象的维度，再定义后文计算预测值和误差时需要的运算。不同单位的特征可以组成向量，但直接计算长度或距离会受单位选择影响。用于距离比较前，通常需要按训练数据确定尺度变换。

\subsection{对象与记号}
标量是单个数。本章除特别说明外使用实数，实数集合记为$\R$。向量默认为列向量，写作$\bm x=(x_1,\ldots,x_d)^{\mathsf T}\in\R^d$。矩阵$A=(A_{ij})\in\R^{m\times d}$有$m$行、$d$列，$A_{ij}$是第$i$行第$j$列的元素。上标$\mathsf T$表示转置，即$(A^{\mathsf T})_{ji}=A_{ij}$，转置后矩阵的维度为$d\times m$。

零向量与零矩阵分别记为$\bm0$与$0$，维度由上下文确定。$I_d\in\R^{d\times d}$表示单位矩阵。$\diag(a_1,\ldots,a_d)$表示对角元素为$a_1,\ldots,a_d$、其余元素为零的对角矩阵。$\mathbb{1}\{B\}$是指示函数，命题$B$成立时为$1$，否则为$0$。本章的$\log$均表示自然对数。

向量或矩阵只有在维度相同时才能逐元素相加。数乘将每个元素乘以同一个标量。若$\bm v_j\in\R^d$、$a_j\in\R$，则$\sum_{j=1}^k a_j\bm v_j\in\R^d$称为线性组合。

\subsection{内积、范数与距离}
实向量的欧氏内积与欧氏范数定义为
\begin{equation}
 \langle\bm x,\bm y\rangle=\bm x^{\mathsf T}\bm y=\sum_{j=1}^d x_jy_j,
 \qquad \|\bm x\|_2=\sqrt{\bm x^{\mathsf T}\bm x}.
\end{equation}
这里$\bm x,\bm y\in\R^d$，内积是标量。内积为零称为正交。范数满足非负性、绝对齐次性和三角不等式，且仅零向量的范数为零。柯西--施瓦茨不等式（Cauchy--Schwarz inequality）为$|\bm x^{\mathsf T}\bm y|\leq\|\bm x\|_2\|\bm y\|_2$，等号成立当且仅当两向量线性相关。

对$1\leq p<\infty$定义$\|\bm x\|_p=(\sum_j|x_j|^p)^{1/p}$，另定义$\|\bm x\|_\infty=\max_j|x_j|$。相应距离为
\begin{align}
 d_1(x,y)&=\|x-y\|_1=\sum_{j=1}^d|x_j-y_j|,\label{eq:math-manhattan-distance}\\
 d_2(x,y)&=\|x-y\|_2=\sqrt{\sum_{j=1}^d(x_j-y_j)^2},\label{eq:math-euclidean-distance}\\
 d_p(x,y)&=\|x-y\|_p=\left(\sum_{j=1}^d|x_j-y_j|^p\right)^{1/p},\quad p\geq1,\label{eq:math-minkowski-distance}\\
 d_\infty(x,y)&=\|x-y\|_\infty=\max_j|x_j-y_j|.\label{eq:math-chebyshev-distance}
\end{align}
这里沿用$x,y\in\R^d$表示两个向量。四式分别称为曼哈顿距离（Manhattan distance）、欧氏距离（Euclidean distance）、闵可夫斯基距离（Minkowski distance）和切比雪夫距离（Chebyshev distance）。度量要求非负、对称、满足三角不等式，且距离为零当且仅当两对象相同。上述距离均为度量。$0<p<1$时同一幂和表达式一般不满足范数的三角不等式。

非零向量的余弦相似度（cosine similarity）与余弦距离（cosine distance）为
\begin{equation}
 s_{\cos}(x,y)=\frac{x^{\mathsf T}y}{\|x\|_2\|y\|_2},\qquad
 d_{\cos}(x,y)=1-s_{\cos}(x,y),\quad x\ne0,\ y\ne0.
 \label{eq:math-cosine-distance}
\end{equation}
相似度等于夹角余弦，范围为$[-1,1]$。余弦距离的范围为$[0,2]$，但一般不是度量，正比例向量可有零距离，三角不等式也未必成立。对等长离散符号序列，汉明距离（Hamming distance）为
\begin{equation}
 d_{\mathrm H}(x,y)=\sum_{j=1}^d\mathbb{1}\{x_j\ne y_j\}.
 \label{eq:math-hamming-distance}
\end{equation}
它计数不同坐标的位置，在固定长度的同一符号集合上满足度量条件。

\subsection{矩阵乘法与线性变换}
若$A\in\R^{m\times d}$、$B\in\R^{d\times k}$，则
\begin{equation}
 (AB)_{i\ell}=\sum_{j=1}^d A_{ij}B_{j\ell},\qquad AB\in\R^{m\times k}.
\end{equation}
相邻的两个维度必须相等。特别地，$A\bm x\in\R^m$且$(A\bm x)_i=\sum_jA_{ij}x_j$。映射$\bm x\mapsto A\bm x$保持加法与数乘，称为线性变换。矩阵乘法满足结合律和分配律，一般不满足交换律。转置满足$(AB)^{\mathsf T}=B^{\mathsf T}A^{\mathsf T}$。对$\bm x\in\R^m$、$\bm y\in\R^d$，外积$\bm x\bm y^{\mathsf T}\in\R^{m\times d}$的元素为$x_iy_j$。

弗罗贝尼乌斯范数（Frobenius norm）与谱范数（spectral norm）定义为
\begin{equation}
 \|A\|_{\mathrm F}=\left(\sum_{i,j}A_{ij}^2\right)^{1/2},\qquad
 \|A\|_2=\sup_{\|\bm x\|_2=1}\|A\bm x\|_2.
\end{equation}
前者计入全部元素，后者度量单位输入向量的最大输出长度。维度相容时$\|AB\|_2\leq\|A\|_2\|B\|_2$。

\section{线性空间与矩阵运算性质}
\subsection{线性无关与秩}
向量组$\bm v_1,\ldots,\bm v_k$线性无关，是指$\sum_j a_j\bm v_j=\bm0$仅在全部$a_j=0$时成立。全部线性组合构成它们张成的空间。子空间的一组线性无关生成向量称为基，基中向量的个数称为维数。

对$A\in\R^{m\times d}$，列空间与零空间为
\begin{equation}
 \mathcal C(A)=\{A\bm x:\bm x\in\R^d\}\subseteq\R^m,\qquad
 \mathcal N(A)=\{\bm x\in\R^d:A\bm x=\bm0\}\subseteq\R^d.
\end{equation}
列空间包含全部可能输出，零空间包含映射为零的输入。行空间为$\mathcal C(A^{\mathsf T})$。秩（rank）定义为列空间的维数，满足
\begin{equation}
 \operatorname{rank}(A)=\operatorname{rank}(A^{\mathsf T}),\qquad
 \operatorname{rank}(A)+\dim\mathcal N(A)=d.
\end{equation}
第二式称为秩--零度定理。正交直和分解为$\R^d=\mathcal C(A^{\mathsf T})\oplus\mathcal N(A)$、$\R^m=\mathcal C(A)\oplus\mathcal N(A^{\mathsf T})$，其中$\oplus$表示两子空间正交，且每个向量有唯一的分解。

秩不超过$\min(m,d)$。秩为$d$称为满列秩，要求$m\geq d$。秩为$m$称为满行秩，要求$m\leq d$。维度相容时，
\begin{equation}
 \operatorname{rank}(AB)\leq\min\{\operatorname{rank}(A),\operatorname{rank}(B)\},\qquad
 \operatorname{rank}(A^{\mathsf T}A)=\operatorname{rank}(A).
\end{equation}
后一结论来自$\mathcal N(A^{\mathsf T}A)=\mathcal N(A)$：$A^{\mathsf T}A\bm x=\bm0$蕴含$\|A\bm x\|_2^2=0$，反向蕴含直接成立。

\subsection{迹与行列式}
方阵$A\in\R^{d\times d}$的迹（trace）为$\operatorname{tr}(A)=\sum_jA_{jj}$。迹是标量，具有线性性质，且$\operatorname{tr}(A^{\mathsf T})=\operatorname{tr}(A)$。若$A\in\R^{m\times d}$、$B\in\R^{d\times m}$，则
\begin{equation}
 \operatorname{tr}(AB)=\sum_{i=1}^m\sum_{j=1}^dA_{ij}B_{ji}=\operatorname{tr}(BA).
\end{equation}
循环移动可推广到多个因子，只要各乘积有定义且是方阵。一般不能任意交换因子。另有$\|A\|_{\mathrm F}^2=\operatorname{tr}(A^{\mathsf T}A)$。若$B\in\R^{d\times d}$，则$\bm x^{\mathsf T}B\bm x=\operatorname{tr}(B\bm x\bm x^{\mathsf T})$。

行列式（determinant）是关于方阵各列交替多线性的标量函数，由$\det I_d=1$确定。绝对值表示体积缩放因子，符号表示定向。同阶方阵满足$\det(AB)=\det A\det B$、$\det(A^{\mathsf T})=\det A$。三角矩阵的行列式为对角元素之积。$|A|$也表示行列式。

\subsection{逆与线性方程组}
若方阵$A\in\R^{d\times d}$存在$A^{-1}$满足$AA^{-1}=A^{-1}A=I_d$，则称为可逆，$A^{-1}$是唯一的逆矩阵（inverse matrix）。可逆等价于$\operatorname{rank}(A)=d$、$\mathcal N(A)=\{\bm0\}$、$\det A\ne0$，也等价于对每个$\bm b\in\R^d$，$A\bm x=\bm b$恰有一个解。可逆同阶方阵满足
\begin{equation}
 (AB)^{-1}=B^{-1}A^{-1},\qquad (A^{\mathsf T})^{-1}=(A^{-1})^{\mathsf T}.
\end{equation}
对一般$A\in\R^{m\times d}$，方程有解当且仅当$\bm b\in\mathcal C(A)$。若$\bm x_0$为一个解，则全部解为$\bm x_0+\mathcal N(A)$。解存在时，唯一性等价于满列秩。矩形矩阵没有上述双侧逆，但满列秩矩阵有左逆，满行秩矩阵有右逆。

\subsection{穆尔--彭罗斯伪逆}
每个$A\in\R^{m\times d}$均有唯一的穆尔--彭罗斯伪逆（Moore--Penrose pseudoinverse）$A^+\in\R^{d\times m}$，定义条件为
\begin{align}
 AA^+A&=A,&A^+AA^+&=A^+,\\
 (AA^+)^{\mathsf T}&=AA^+,&(A^+A)^{\mathsf T}&=A^+A.
\end{align}
可逆方阵的伪逆等于逆。满列秩时$A^+=(A^{\mathsf T}A)^{-1}A^{\mathsf T}$，满行秩时$A^+=A^{\mathsf T}(AA^{\mathsf T})^{-1}$。一般构造将在奇异值分解中给出。

$AA^+$是到列空间的正交投影，$A^+A$是到行空间的正交投影。对任意$\bm b\in\R^m$，$A^+\bm b$是使$\|A\bm x-\bm b\|_2$最小的全部解中范数最小的唯一解。全部最小平方残差解为
\begin{equation}
 \bm x=A^+\bm b+(I_d-A^+A)\bm z,\qquad\bm z\in\R^d.
\end{equation}
其输出均为$\bm b$到列空间的投影。第二项属于零空间，与行空间中的$A^+\bm b$正交，加入非零第二项会增大解的范数。

\section{正交投影与正定性}
\subsection{正交矩阵与投影}
方阵$U\in\R^{d\times d}$满足$U^{\mathsf T}U=I_d$时称为正交矩阵（orthogonal matrix）。它也满足$UU^{\mathsf T}=I_d$，故$U^{-1}=U^{\mathsf T}$。正交变换保持内积与欧氏范数。

若$U_r\in\R^{d\times r}$满足$U_r^{\mathsf T}U_r=I_r$，其列为一个$r$维子空间的标准正交基。令$P=U_rU_r^{\mathsf T}$，则
\begin{equation}
 \bm x=P\bm x+(I_d-P)\bm x,\qquad
 \|\bm x\|_2^2=\|P\bm x\|_2^2+\|(I_d-P)\bm x\|_2^2.
\end{equation}
两项分别位于该子空间及其正交补。投影矩阵满足$P=P^{\mathsf T}$和$P^2=P$。幂等性表示重复投影不改变结果，对称性进一步保证投影方向与目标子空间正交。对称且幂等的实矩阵均为正交投影矩阵。

图\ref{fig:math-orthogonal-projection}取$\bm x=(2,1)^{\mathsf T}$，并投影到第一个坐标轴。投影保留第一坐标，残差保留第二坐标。两者正交，因此长度平方满足$5=4+1$。后文最小二乘会把同样的分解用于样本预测向量。
\begin{figure}[htbp]
\centering
\begin{tikzpicture}[>=Latex,scale=1.3,every node/.style={font=\small}]
 \draw[->] (-0.3,0)--(3,0) node[right] {$x_1$};
 \draw[->] (0,-0.2)--(0,1.7) node[above] {$x_2$};
 \draw[->,thick] (0,0)--(2,1) node[above right] {$\bm x=(2,1)^{\mathsf T}$};
 \draw[->,thick,blue] (0,0)--(2,0) node[below] {$P\bm x=(2,0)^{\mathsf T}$};
 \draw[->,thick,dashed] (2,0)--(2,1) node[midway,right] {$\bm x-P\bm x$};
 \draw (1.85,0)--(1.85,0.15)--(2,0.15);
\end{tikzpicture}
\caption{向量到横轴的正交投影。实线箭头分别表示原向量和投影向量，虚线箭头表示与横轴垂直的残差。}
\label{fig:math-orthogonal-projection}
\end{figure}

单位向量$\bm u$张成的直线具有投影$\bm u\bm u^{\mathsf T}\bm x$。其系数由
\begin{equation}
 \frac{d}{da}\|\bm x-a\bm u\|_2^2=-2\bm u^{\mathsf T}\bm x+2a=0
 \label{eq:math-projection-stationarity}
\end{equation}
确定。这里$\|\bm u\|_2=1$、$a\in\R$，二阶导数为$2$，故$a=\bm u^{\mathsf T}\bm x$是唯一最小点。非单位的非零向量具有系数$\bm u^{\mathsf T}\bm x/(\bm u^{\mathsf T}\bm u)$。

\subsection{二次型与正定矩阵}
实方阵$A\in\R^{d\times d}$满足$A=A^{\mathsf T}$时称为对称矩阵（symmetric matrix）。本章对实对称矩阵定义正定性：
\begin{align}
 A\succeq0&\ \Longleftrightarrow\ \bm v^{\mathsf T}A\bm v\geq0\quad\text{对所有 }\bm v\in\R^d,\\
 A\succ0&\ \Longleftrightarrow\ \bm v^{\mathsf T}A\bm v>0\quad\text{对所有非零 }\bm v\in\R^d.
\end{align}
第一种性质称为半正定（positive semidefinite，PSD），第二种称为正定（positive definite，PD）。标量$\bm v^{\mathsf T}A\bm v$称为二次型。负定与半负定分别由$-A\succ0$与$-A\succeq0$定义。二次型既取正值又取负值时称为不定。任意实方阵的二次型只取决于对称部分$(A+A^{\mathsf T})/2$。

对称矩阵之间记$A\succeq B$表示$A-B\succeq0$。半正定矩阵的非负系数线性组合仍半正定。若$A\succeq0$、$C\in\R^{d\times k}$，则$C^{\mathsf T}AC\succeq0$，因为对$\bm z\in\R^k$有
\begin{equation}
 \bm z^{\mathsf T}C^{\mathsf T}AC\bm z=(C\bm z)^{\mathsf T}A(C\bm z)\geq0.
\end{equation}
若$A\succ0$，则$C^{\mathsf T}AC$正定当且仅当$C$满列秩。任意$B\in\R^{m\times d}$都有$B^{\mathsf T}B\succeq0$，且正定当且仅当$B$满列秩。

正定矩阵必可逆，逆矩阵也正定。半正定矩阵可能奇异，但任意$\epsilon>0$均使$A+\epsilon I_d\succ0$。西尔维斯特判据（Sylvester criterion）指出，对称矩阵正定等价于全部顺序主子式为正。顺序主子式是前$k$行、前$k$列形成方阵的行列式。半正定等价于全部主子式非负，主子式允许选取任意相同的行列指标，仅检查顺序主子式不足以判定半正定。

\subsection{分块矩阵与舒尔补}
设对称分块矩阵为
\begin{equation}
 M=\begin{pmatrix}A&B\\B^{\mathsf T}&C\end{pmatrix},\qquad
 A\in\R^{d\times d},\ B\in\R^{d\times k},\ C\in\R^{k\times k}.
\end{equation}
当$A$可逆时，$S=C-B^{\mathsf T}A^{-1}B\in\R^{k\times k}$称为关于$A$的舒尔补（Schur complement）。消去分块方程中的第一组未知量后，第二组未知量所满足的方程就由$S$决定。下面进一步假设$A\succ0$，用配方判断整个矩阵的正定性。配方给出
\begin{equation}
 \begin{pmatrix}\bm x\\\bm y\end{pmatrix}^{\mathsf T}M
 \begin{pmatrix}\bm x\\\bm y\end{pmatrix}
 =(\bm x+A^{-1}B\bm y)^{\mathsf T}A(\bm x+A^{-1}B\bm y)+\bm y^{\mathsf T}S\bm y.
\end{equation}
这里$\bm x\in\R^d$、$\bm y\in\R^k$。在$A\succ0$的前提下，$M\succ0$等价于$S\succ0$，$M\succeq0$等价于$S\succeq0$。

\section{矩阵分解}
矩阵分解将一个矩阵表示为结构较简单的矩阵的乘积。三角矩阵便于逐项解方程，正交矩阵保持欧氏长度，对角矩阵使各坐标的作用分离。各分解的存在条件不同。

\subsection{LU 分解}
下三角--上三角分解（Lower--Upper decomposition，LU 分解）写为$A=LU$，其中$L$为单位下三角矩阵，$U$为上三角矩阵。非奇异方阵$A\in\R^{d\times d}$的全部顺序主子式非零时，这样的分解存在且唯一。

一般非奇异方阵需要允许交换行，形式为
\begin{equation}
 PA=LU,\qquad P,L,U\in\R^{d\times d}.
\end{equation}
$P$为置换矩阵，满足$P^{-1}=P^{\mathsf T}$。通过选择非零主元，任意非奇异方阵均有这一分解。求解$A\bm x=\bm b$可依次解$L\bm z=P\bm b$、$U\bm x=\bm z$，分别从首行向末行、从末行向首行代入。选主元有助于控制舍入误差，但不能消除方程对扰动的敏感性。

\subsection{QR 分解}
正交--三角分解（Orthogonal--Triangular decomposition，QR 分解）对任意$A\in\R^{m\times d}$存在：
\begin{equation}
 A=QR,\qquad Q\in\R^{m\times m},\quad Q^{\mathsf T}Q=I_m,\quad R\in\R^{m\times d}.
\end{equation}
$R$为上梯形矩阵，即$i>j$时$R_{ij}=0$。当$m\geq d$时，可取简化形式$A=Q_1R_1$，其中$Q_1\in\R^{m\times d}$、$Q_1^{\mathsf T}Q_1=I_d$、$R_1\in\R^{d\times d}$为上三角矩阵。$A$满列秩且规定$R_1$对角元素为正时，简化分解唯一。

正交分解给出
\begin{equation}
 \|A\bm x-\bm b\|_2^2=\|R_1\bm x-Q_1^{\mathsf T}\bm b\|_2^2
 +\|(I_m-Q_1Q_1^{\mathsf T})\bm b\|_2^2.
\end{equation}
这里$\bm x\in\R^d$、$\bm b\in\R^m$。满列秩时，最小平方残差解由$R_1\bm x=Q_1^{\mathsf T}\bm b$唯一确定。秩不足时$R_1$奇异，不能使用唯一解结论。豪斯霍尔德反射（Householder reflection）和吉文斯旋转（Givens rotation）可通过正交变换构造 QR 分解。

\subsection{特征分解}
若$A\in\R^{d\times d}$与非零向量$\bm u$满足$A\bm u=\lambda\bm u$，则$\lambda$称为特征值（eigenvalue），$\bm u$称为特征向量（eigenvector）。特征值是$\det(A-\lambda I_d)=0$的根，一般实矩阵可能有复特征值。

在所讨论数域上，$A$有$d$个线性无关的特征向量时称为可对角化，此时
\begin{equation}
 A=V\Lambda V^{-1},\qquad\Lambda=\diag(\lambda_1,\ldots,\lambda_d).
\end{equation}
$V$的列为特征向量。一般方阵未必可对角化。不同特征值的特征向量线性无关，因此$d$个互异特征值足以保证可对角化。

实对称矩阵的谱定理（spectral theorem）保证特征值为实数，且
\begin{equation}
 A=U\Lambda U^{\mathsf T},\qquad U^{\mathsf T}U=I_d.
\end{equation}
由$\bm v=U\bm z$得到$\bm v^{\mathsf T}A\bm v=\sum_j\lambda_jz_j^2$。所以半正定等价于全部特征值非负，正定等价于全部特征值为正。另有
\begin{equation}
 \operatorname{tr}(A)=\sum_j\lambda_j,\qquad\det A=\prod_j\lambda_j,\qquad
 \lambda_{\min}\leq\frac{\bm v^{\mathsf T}A\bm v}{\bm v^{\mathsf T}\bm v}\leq\lambda_{\max}.
\end{equation}
最后一式要求$\bm v\ne\bm0$，比值称为瑞利商（Rayleigh quotient），端点由对应特征向量达到。一般方阵的迹与行列式仍等于特征值之和与之积，按代数重数计入。

半正定矩阵有唯一的对称半正定平方根$A^{1/2}=U\diag(\sqrt{\lambda_j})U^{\mathsf T}$。正定时还可定义$A^{-1/2}=U\diag(\lambda_j^{-1/2})U^{\mathsf T}$。重特征值空间内的标准正交基不唯一，但这些矩阵函数不依赖基的选择。

\subsection{奇异值分解}
奇异值分解（Singular Value Decomposition，SVD）对任意$A\in\R^{m\times d}$存在：
\begin{equation}
 A=U\Sigma V^{\mathsf T},\qquad U\in\R^{m\times m},\quad
 V\in\R^{d\times d},\quad\Sigma\in\R^{m\times d}.
\end{equation}
$U,V$均正交。令$q=\min(m,d)$，$\Sigma$只在$(j,j)$位置可能非零，且$\sigma_1\geq\cdots\geq\sigma_q\geq0$，称为奇异值（singular values）。此处$\Sigma$为奇异值矩阵，概率部分的同一字母表示协方差矩阵，须按所在公式区分。

若$r=\operatorname{rank}(A)$，则恰有$r$个正奇异值。取相应列组成$U_r\in\R^{m\times r}$、$V_r\in\R^{d\times r}$，得到
\begin{align}
 A&=U_r\diag(\sigma_1,\ldots,\sigma_r)V_r^{\mathsf T},\\
 A^+&=V_r\diag(\sigma_1^{-1},\ldots,\sigma_r^{-1})U_r^{\mathsf T}.
\end{align}
零矩阵的伪逆为零。也可写$A^+=V\Sigma^+U^{\mathsf T}$，$\Sigma^+\in\R^{d\times m}$将正奇异值取倒数，将零奇异值保持为零。

因为$A^{\mathsf T}A=V\Sigma^{\mathsf T}\Sigma V^{\mathsf T}$，其非零特征值为$\sigma_j^2$，且$\|A\|_2=\sigma_1$、$\|A\|_{\mathrm F}^2=\sum_j\sigma_j^2$。对$0\leq k<r$，截断矩阵$A_k=\sum_{j=1}^k\sigma_j\bm u_j\bm v_j^{\mathsf T}$的秩至多为$k$，满足低秩近似定理：
\begin{equation}
 \min_{\operatorname{rank}(B)\leq k}\|A-B\|_2=\sigma_{k+1},\qquad
 \min_{\operatorname{rank}(B)\leq k}\|A-B\|_{\mathrm F}^2=\sum_{j=k+1}^q\sigma_j^2.
\end{equation}
$A_k$达到这两个最小值，其中$\bm u_j$和$\bm v_j$分别为$U$和$V$的第$j$列，$k$为整数。把正的小奇异值置零后再取倒数，得到截断矩阵的伪逆，一般不同于原矩阵的精确伪逆。截断舍去数据约束较弱的方向，可能减小噪声影响，也会改变原来的最小二乘解。

\subsection{Cholesky 与对称分解}
实对称正定矩阵$A\in\R^{d\times d}$具有唯一的乔列斯基分解（Cholesky decomposition）$A=LL^{\mathsf T}$，其中$L$是对角元素严格为正的下三角矩阵。反过来，可逆下三角矩阵$L$使$LL^{\mathsf T}$正定。解方程可依次解$L\bm z=\bm b$与$L^{\mathsf T}\bm x=\bm z$。同时$\det A=\prod_jL_{jj}^2$，故$\log\det A=2\sum_j\log L_{jj}$。

半正定矩阵可写为$BB^{\mathsf T}$，但$B$可能秩不足，不能要求严格正对角元素的 Cholesky 因子。对称不定矩阵可使用带对称置换的$PAP^{\mathsf T}=LDL^{\mathsf T}$，其中$L$为单位下三角矩阵，$D$为具有$1\times1$或$2\times2$对角块的对称分块对角矩阵。允许$2\times2$块可处理无法使用标量主元的情形。

\subsection{条件数与投影恒等式}
可逆方阵的二范数条件数（condition number）为
\begin{equation}
 \kappa_2(A)=\|A\|_2\|A^{-1}\|_2=\frac{\sigma_{\max}(A)}{\sigma_{\min}(A)}\geq1.
\end{equation}
若$A\bm x=\bm b$且只扰动右端，$\delta\bm x=A^{-1}\delta\bm b$。当$\bm b\ne\bm0$时，
\begin{equation}
 \frac{\|\delta\bm x\|_2}{\|\bm x\|_2}
 \leq\kappa_2(A)\frac{\|\delta\bm b\|_2}{\|\bm b\|_2}.
\end{equation}
可逆性保证解唯一，条件数描述对扰动的敏感性。条件数很大时称为病态。满列秩矩形矩阵也可定义$\kappa_2(A)=\sigma_1/\sigma_d$，此时$\kappa_2(A^{\mathsf T}A)=\kappa_2(A)^2$。数学秩以严格为零为准，数值秩还需给定尺度和容差。

考虑用一个低维线性子空间重建一组向量。设$\bm z_i\in\R^d$，$i=1,\ldots,n$且$n\geq1$。令$S=n^{-1}\sum_{i=1}^n\bm z_i\bm z_i^{\mathsf T}$，并令$U_r\in\R^{d\times r}$满足$U_r^{\mathsf T}U_r=I_r$，其中$1\leq r\leq d$。矩阵$U_rU_r^{\mathsf T}$把每个向量投影到它的列空间。由正交分解与迹的性质，
\begin{align}
 \frac1n\sum_i\|\bm z_i-U_rU_r^{\mathsf T}\bm z_i\|_2^2
 &=\frac1n\sum_i(\|\bm z_i\|_2^2-\|U_r^{\mathsf T}\bm z_i\|_2^2)\\
 &=\operatorname{tr}(S)-\operatorname{tr}(U_r^{\mathsf T}SU_r).
 \label{eq:math-pca-reconstruction}
\end{align}
这里$S\in\R^{d\times d}$半正定，不要求$\bm z_i$来自随机抽样。若$\bm z_i$是减去样本均值后的向量，则$S$是以$n$为分母的样本协方差矩阵。未中心化时，它是二阶矩矩阵，最优子空间被限制为经过原点。若$S$的特征值按$\lambda_1\geq\cdots\geq\lambda_d\geq0$排列，则在全部$r$维子空间上，右端最小值为$\sum_{j>r}\lambda_j$。令$\bm u_j$为$S$的第$j$个单位特征向量，则保留的迹为$\sum_j\lambda_jw_j$，其中$w_j=\|U_r^{\mathsf T}\bm u_j\|_2^2$表示该特征方向保留的平方投影长度。正交投影保证$0\leq w_j\leq1$且$\sum_jw_j=r$。把权重配置给最大的$r$个特征值可使该迹达到最大值，选择相应特征向量作为$U_r$的列即可实现。若截断位置的特征值相同，最优子空间可能不唯一。

\section{微积分工具}
\subsection{导数、全微分与链式法则}
函数$f$在$x$的导数为$f'(x)=\lim_{h\to0}[f(x+h)-f(x)]/h$，前提是极限存在。对$f:\R^d\to\R$，偏导数只沿一个坐标方向取极限。各偏导数存在并不自动保证全微分存在。若存在$\bm g\in\R^d$使
\begin{equation}
 f(\bm x+\bm h)=f(\bm x)+\bm g^{\mathsf T}\bm h+o(\|\bm h\|_2),
\end{equation}
则$f$在$\bm x$可微。$\bm g=\nabla f(\bm x)$称为梯度（gradient），分量为各偏导数。小量$o(\|\bm h\|_2)$除以$\|\bm h\|_2$后趋于零。偏导数在邻域内连续是可微的充分条件。

对$\bm f:\R^d\to\R^m$，雅可比矩阵（Jacobian matrix）为$D\bm f=(\partial f_i/\partial x_j)_{i,j}\in\R^{m\times d}$。可微时，$\bm f(\bm x+\bm h)=\bm f(\bm x)+D\bm f(\bm x)\bm h+o(\|\bm h\|_2)$。若$\bm g:\R^m\to\R^k$也可微，链式法则（chain rule）给出
\begin{equation}
 D(\bm g\circ\bm f)(\bm x)=D\bm g(\bm f(\bm x))D\bm f(\bm x)\in\R^{k\times d}.
\end{equation}
当$g$为标量函数时，$\nabla(g\circ\bm f)=(D\bm f)^{\mathsf T}\nabla g\in\R^d$。

\subsection{二阶导数与泰勒展开}
海森矩阵（Hessian matrix）为$H_f=(\partial^2 f/\partial x_i\partial x_j)_{i,j}\in\R^{d\times d}$。若二阶偏导数在邻域内连续，混合偏导数可交换，$H_f$对称，且二阶泰勒展开（Taylor expansion）为
\begin{equation}
 f(\bm x+\bm h)=f(\bm x)+\nabla f(\bm x)^{\mathsf T}\bm h
 +\frac12\bm h^{\mathsf T}H_f(\bm x)\bm h+o(\|\bm h\|_2^2).
\end{equation}
可微函数沿方向$\bm v$的方向导数为$\nabla f^{\mathsf T}\bm v$。梯度非零时，柯西--施瓦茨不等式说明单位方向中的最小方向导数为$-\|\nabla f\|_2$，由负梯度的单位方向取得。这保证足够小的正位移使函数下降。

定义域内部的可微局部极小点必须满足$\nabla f=\bm0$。二阶连续可微时，该点还必须满足$H_f\succeq0$。若驻点满足$H_f\succ0$，则为严格局部极小点。仅半正定不足以判定极小点。

\subsection{矩阵微分}
对固定$A\in\R^{d\times d}$、$\bm b\in\R^d$，
\begin{equation}
 \nabla_{\bm x}(\bm b^{\mathsf T}\bm x)=\bm b,\qquad
 \nabla_{\bm x}(\bm x^{\mathsf T}A\bm x)=(A+A^{\mathsf T})\bm x.
\end{equation}
第二式由$d(\bm x^{\mathsf T}A\bm x)=(d\bm x)^{\mathsf T}A\bm x+\bm x^{\mathsf T}A\,d\bm x$得到。$A$对称时梯度为$2A\bm x$。

对固定$A\in\R^{m\times d}$、$\bm b\in\R^m$，$F(\bm x)=\tfrac12\|A\bm x-\bm b\|_2^2$满足
\begin{equation}
 dF=(A\bm x-\bm b)^{\mathsf T}A\,d\bm x,\quad
 \nabla F=A^{\mathsf T}(A\bm x-\bm b),\quad H_F=A^{\mathsf T}A.
\end{equation}
梯度为$d$维，海森矩阵为$d\times d$且半正定。

对可微变化的可逆方阵，
\begin{equation}
 d(A^{-1})=-A^{-1}(dA)A^{-1},\qquad
 d\log|\det A|=\operatorname{tr}(A^{-1}dA).
\end{equation}
第一式由$d(AA^{-1})=0$推得，第二式在可逆矩阵邻域内成立。正定时可写$\log\det A$。对矩阵函数$F(A)$，可由$dF=\operatorname{tr}((\nabla_A F)^{\mathsf T}dA)$定义同维度的矩阵梯度。

\subsection{积分与变量替换}
若$f$连续且$F'=f$，微积分基本定理给出$\int_a^b f(x)\,dx=F(b)-F(a)$。若$u,v$连续可微，分部积分为$\int_a^b uv'\,dx=[uv]_a^b-\int_a^b u'v\,dx$。

若$\bm x=\bm g(\bm z)$为开集之间一一、连续可微且逆也连续可微的变换，则
\begin{equation}
 \int_{\mathcal X}f(\bm x)\,d\bm x
 =\int_{\mathcal Z}f(\bm g(\bm z))|\det D\bm g(\bm z)|\,d\bm z.
\end{equation}
这里$\mathcal X,\mathcal Z\subseteq\R^d$，雅可比矩阵为$d\times d$，要求$f$非负或绝对可积。行列式绝对值是体积因子。非负函数可交换积分次序，这是托内利定理（Tonelli theorem）。绝对可积函数也可交换次序，这是富比尼定理（Fubini theorem）。缺少这些条件时，不能任意交换积分或求和次序。

\section{概率空间、随机变量与分布}
即使输入读数相同，设备下一小时的温度也可能不同。概率模型需要区分可能发生的结果、结果组成的事件，以及事件的概率。先建立这三个对象，才能准确解释条件概率和预测分布。下面的抽象定义统一覆盖离散结果与连续测量，后续计算主要使用质量函数和密度。

\subsection{事件与概率}
概率空间（probability space）是三元组$(\Omega,\mathcal F,P)$。$\Omega$为所有可能结果组成的样本空间。$\mathcal F$为事件组成的$\sigma$代数，包含$\Omega$，对补集与可数并封闭。概率测度$P$定义在$\mathcal F$上，满足
\begin{equation}
 P(A)\geq0,\qquad P(\Omega)=1,\qquad
 P\left(\bigcup_{i=1}^\infty A_i\right)=\sum_{i=1}^\infty P(A_i)
\end{equation}
最后一式要求$A_i$两两不交。公理蕴含$P(\varnothing)=0$、$P(A^{\mathrm c})=1-P(A)$及$P(A\cup B)=P(A)+P(B)-P(A\cap B)$。

若$P(B)>0$，条件概率（conditional probability）定义为$P(A\mid B)=P(A\cap B)/P(B)$。它在$B$内重新归一化。事件$A,B$独立（independent）是指$P(A\cap B)=P(A)P(B)$。相互独立要求任意有限子集的交事件概率等于相应概率之积，比两两独立更强。

\subsection{随机变量与分布函数}
随机变量（random variable）是可测映射$X:\Omega\to\R$，即对每个实数$t$都有$\{X\leq t\}\in\mathcal F$。分布由$P_X(B)=P(X\in B)$确定，$B$为实数上的可测集合。大写$X$表示随机变量，小写$x$表示可能取值。分布函数（cumulative distribution function，CDF）为$F_X(t)=P(X\leq t)$，它单调不减、右连续，在$t\to-\infty$与$t\to+\infty$时分别趋于$0$与$1$。区间概率为$P(a<X\leq b)=F_X(b)-F_X(a)$。

随机变量用于表示观测结果的数值化对象。例如，$X$可表示某时刻的温度，$Y$可表示随后一小时的温度。离散随机变量的取值至多可数，概率质量函数（probability mass function，PMF）为$p_X(x)=P(X=x)$，满足$p_X(x)\geq0$与$\sum_xp_X(x)=1$。分布相对于勒贝格测度绝对连续时，存在概率密度函数（probability density function，PDF），满足
\begin{equation}
 P(X\in B)=\int_Bp_X(x)\,dx,\qquad p_X(x)\geq0,\qquad\int_{\R}p_X(x)\,dx=1.
\end{equation}
这时$F_X(t)=\int_{-\infty}^t p_X(x)\,dx$，几乎处处有$F_X'=p_X$。密度可大于$1$，但单点概率为零。密度量纲是变量量纲的倒数，积分后概率无量纲。一般分布可能有离散或奇异连续部分，不能假定每个分布都有普通密度。

随机向量$\bm X=(X_1,\ldots,X_d)^{\mathsf T}$是到$\R^d$的可测映射。多维密度相对于$d$维勒贝格测度定义，积分覆盖全部坐标。

\subsection{联合、边缘、条件与独立}
联合分布（joint distribution）描述共同取值。有联合密度$p(x,y)$时，边缘密度（marginal density）为$p_X(x)=\int p(x,y)\,dy$。离散变量使用求和。无歧义时用$p(x)$代替$p_X(x)$。

在$p(x)>0$处，条件密度或条件质量函数满足
\begin{equation}
 p(x,y)=p(y\mid x)p(x),\qquad p(y\mid x)=\frac{p(x,y)}{p(x)}.
 \label{eq:math-conditional-probability}
\end{equation}
条件分布在$y$上归一化。连续情形不能将其写成两个单点概率相除。条件分布仅需在$X$的分布下几乎处处定义，$p(x)=0$处不由比值确定。多变量链式法则为$p(x_1,\ldots,x_d)=p(x_1)\prod_{j=2}^d p(x_j\mid x_1,\ldots,x_{j-1})$，在条件分布有定义处使用。

$X,Y$独立，是指任意可测集合$B,C$满足$P(X\in B,Y\in C)=P(X\in B)P(Y\in C)$。有联合密度时，等价于$p(x,y)=p(x)p(y)$几乎处处成立。给定$Z$时条件独立（conditional independence）记为$X\perp Y\mid Z$，密度形式为
\begin{equation}
 p(x,y\mid z)=p(x\mid z)p(y\mid z)
\end{equation}
对$Z$的分布下几乎所有$z$成立。独立与条件独立一般不能互相推出。

\section{期望、方差与协方差}
\subsection{期望与矩}
若$\E|g(X)|<\infty$，期望（expectation）为
\begin{equation}
 \E[g(X)]=\sum_xg(x)p_X(x)\quad\text{或}\quad
 \E[g(X)]=\int g(x)p_X(x)\,dx.
\end{equation}
前式针对离散分布，后式针对有密度的分布。一般定义是对分布测度积分。非负函数允许期望为$+\infty$，但正负部分均为无穷时不能定义为它们的差。均值为$\mu=\E X$，$k$阶原点矩为$\E[X^k]$，中心矩为$\E[(X-\mu)^k]$，均需检查存在条件。

期望具有线性性质，不要求变量独立。若$X,Y$独立且绝对可积，则$\E[XY]=\E X\E Y$。对凸函数$\phi$，若$X$与$\phi(X)$可积，詹森不等式（Jensen inequality）为$\phi(\E X)\leq\E[\phi(X)]$。

\subsection{方差与协方差}
若$\E[X^2]<\infty$，方差（variance）与标准差（standard deviation）为
\begin{equation}
 \Var(X)=\E[(X-\mu)^2]=\E[X^2]-\mu^2,\qquad
 \operatorname{sd}(X)=\sqrt{\Var(X)}.
\end{equation}
方差非负，且为零当且仅当$X=\mu$几乎必然成立。方差量纲为$X$量纲的平方，标准差与$X$同量纲。常数$a,b$满足$\Var(aX+b)=a^2\Var(X)$。

若$X,Y$有有限二阶矩，协方差（covariance）为
\begin{align}
 \operatorname{Cov}(X,Y)&=\E[(X-\E X)(Y-\E Y)],\\
 \Var(X+Y)&=\Var(X)+\Var(Y)+2\operatorname{Cov}(X,Y).
\end{align}
独立蕴含协方差为零，反向结论一般不成立。两方差均为正时，相关系数（correlation coefficient）为$\rho_{XY}=\operatorname{Cov}(X,Y)/(\operatorname{sd}(X)\operatorname{sd}(Y))\in[-1,1]$，它无量纲。任一方差为零时，相关系数无定义。

有有限二阶矩的$\bm X\in\R^d$具有均值$\bm\mu=\E\bm X$及协方差矩阵$\Sigma=\E[(\bm X-\bm\mu)(\bm X-\bm\mu)^{\mathsf T}]\in\R^{d\times d}$。第$(i,j)$个元素为$\operatorname{Cov}(X_i,X_j)$，且
\begin{equation}
 \bm v^{\mathsf T}\Sigma\bm v=\E[(\bm v^{\mathsf T}(\bm X-\bm\mu))^2]
 =\Var(\bm v^{\mathsf T}\bm X)\geq0.
 \label{eq:math-covariance-psd}
\end{equation}
此式对任意$\bm v\in\R^d$成立，证明$\Sigma$对称半正定。它正定当且仅当每个非零线性组合均有正方差。固定$B\in\R^{m\times d}$、$\bm b\in\R^m$时，$\E[B\bm X+\bm b]=B\bm\mu+\bm b$，$\operatorname{Cov}(B\bm X+\bm b)=B\Sigma B^{\mathsf T}$。

\subsection{条件期望与全方差}
对可积$Y$，条件期望$\E[Y\mid X]$是$X$的可积函数，满足对任意有界可测$g$都有$\E[g(X)\E[Y\mid X]]=\E[g(X)Y]$。它几乎处处唯一。存在条件密度时，$\mu(x)=\E[Y\mid X=x]=\int yp(y\mid x)\,dy$，离散时使用求和。

条件期望具有线性性质，满足塔式法则$\E[\E[Y\mid X]]=\E Y$。有界可测$g$还满足$\E[g(X)Y\mid X]=g(X)\E[Y\mid X]$。

若$\E[Y^2]<\infty$，定义$\Var(Y\mid X)=\E[(Y-\E[Y\mid X])^2\mid X]$。对任意常数$a\in\R$，在条件二阶矩有限处展开平方：
\begin{align}
 \E[(Y-a)^2\mid X=x]
 &=\E[(Y-\mu(x))^2\mid X=x]+(\mu(x)-a)^2\notag\\
 &\quad+2(\mu(x)-a)\E[Y-\mu(x)\mid X=x]\notag\\
 &=\Var(Y\mid X=x)+(\mu(x)-a)^2.
 \label{eq:math-conditional-mse}
\end{align}
交叉项由条件均值为零而消失，所以条件平方偏差在$a=\mu(x)$处取得唯一最小值，不需要分布形状假设。

将$Y-\E Y=(Y-\E[Y\mid X])+(\E[Y\mid X]-\E Y)$平方并取期望。交叉项先对$X$条件化后为零，故
\begin{equation}
 \Var(Y)=\E[\Var(Y\mid X)]+\Var(\E[Y\mid X]).
 \label{eq:math-total-variance}
\end{equation}
这称为全方差公式（law of total variance），要求有限二阶矩。

\section{常用概率分布}
本节给出分布的定义、参数范围与基本性质。质量函数与密度在指定支撑集之外均取零。$X\sim\mathcal D$表示$X$服从分布$\mathcal D$。

\subsection{离散分布}
伯努利分布（Bernoulli distribution）$X\sim\operatorname{Bernoulli}(\pi)$定义在$\{0,1\}$上，$0\leq\pi\leq1$，满足$P(X=1)=\pi$、$P(X=0)=1-\pi$。均值为$\pi$，方差为$\pi(1-\pi)$。

二项分布（binomial distribution）$X\sim\operatorname{Binomial}(n,\pi)$的质量函数为
\begin{equation}
 P(X=k)=\binom nk\pi^k(1-\pi)^{n-k},\qquad k=0,\ldots,n.
\end{equation}
$n$为非负整数，$0\leq\pi\leq1$，端点参数按退化分布理解。均值为$n\pi$，方差为$n\pi(1-\pi)$。$n$个独立同参数伯努利变量之和服从该分布。

类别分布（categorical distribution）定义在$\{1,\ldots,K\}$上，$P(X=j)=\pi_j$，$\pi_j\geq0$且$\sum_j\pi_j=1$。多项分布（multinomial distribution）$\bm N\sim\operatorname{Multinomial}(n,\bm\pi)$定义在非负整数向量$\bm k$且$\sum_jk_j=n$上：
\begin{equation}
 P(\bm N=\bm k)=\frac{n!}{\prod_{j=1}^Kk_j!}\prod_{j=1}^K\pi_j^{k_j}.
\end{equation}
$n$为非负整数，$\bm\pi$为类别概率向量，$0^0$按$1$处理。均值为$\E N_j=n\pi_j$，方差为$n\pi_j(1-\pi_j)$，$i\ne j$时协方差为$-n\pi_i\pi_j$。

泊松分布（Poisson distribution）$X\sim\operatorname{Poisson}(\lambda)$为
\begin{equation}
 P(X=k)=e^{-\lambda}\frac{\lambda^k}{k!},\qquad k=0,1,\ldots,\quad\lambda>0.
\end{equation}
均值与方差均为$\lambda$。相互独立的泊松变量之和仍为泊松变量，参数为各参数之和。

\subsection{连续分布}
均匀分布（uniform distribution）$X\sim\operatorname{Uniform}(a,b)$要求$a<b$，在$[a,b]$上的密度为$1/(b-a)$，均值为$(a+b)/2$，方差为$(b-a)^2/12$。

指数分布（exponential distribution）$X\sim\operatorname{Exponential}(\lambda)$的密度为$\lambda e^{-\lambda x}$，$x\geq0$、$\lambda>0$。均值为$1/\lambda$，方差为$1/\lambda^2$。无记忆性质为$P(X>s+t\mid X>s)=P(X>t)$，$s,t\geq0$。

伽马分布（gamma distribution）采用形状--率参数化：
\begin{equation}
 p(x)=\frac{\beta^\alpha}{\Gamma(\alpha)}x^{\alpha-1}e^{-\beta x},
 \qquad x>0,\quad\alpha,\beta>0.
\end{equation}
$\Gamma(\alpha)=\int_0^\infty t^{\alpha-1}e^{-t}\,dt$为伽马函数。均值为$\alpha/\beta$，方差为$\alpha/\beta^2$。若使用尺度参数，尺度为$1/\beta$，须相应改写公式。

贝塔分布（beta distribution）$X\sim\operatorname{Beta}(\alpha,\beta)$的密度为
\begin{equation}
 p(x)=\frac{x^{\alpha-1}(1-x)^{\beta-1}}{\mathrm B(\alpha,\beta)},
 \qquad0<x<1,\quad\alpha,\beta>0.
\end{equation}
$\mathrm B(\alpha,\beta)=\Gamma(\alpha)\Gamma(\beta)/\Gamma(\alpha+\beta)$为贝塔函数。均值为$\alpha/(\alpha+\beta)$，方差为$\alpha\beta/[(\alpha+\beta)^2(\alpha+\beta+1)]$。

狄利克雷分布（Dirichlet distribution）$\bm X\sim\operatorname{Dirichlet}(\alpha_1,\ldots,\alpha_K)$要求$K\geq2$、$\alpha_j>0$。在$x_j>0$、$\sum_jx_j=1$的单纯形上，以$x_1,\ldots,x_{K-1}$为自由坐标，其密度为
\begin{equation}
 p(\bm x)=\frac{\Gamma(\alpha_0)}{\prod_j\Gamma(\alpha_j)}\prod_{j=1}^Kx_j^{\alpha_j-1},
 \qquad\alpha_0=\sum_{j=1}^K\alpha_j.
\end{equation}
密度相对于自由坐标的$(K-1)$维勒贝格测度定义。均值为$\alpha_j/\alpha_0$，方差为$\alpha_j(\alpha_0-\alpha_j)/[\alpha_0^2(\alpha_0+1)]$，$i\ne j$时协方差为$-\alpha_i\alpha_j/[\alpha_0^2(\alpha_0+1)]$。

\subsection{正态分布}
正态分布（normal distribution），又称高斯分布（Gaussian distribution），$X\sim\mathcal N(\mu,\sigma^2)$的密度为
\begin{equation}
 p(x)=\frac1{\sqrt{2\pi}\sigma}\exp\left[-\frac{(x-\mu)^2}{2\sigma^2}\right],
 \qquad x\in\R,\quad\mu\in\R,\quad\sigma>0.
\end{equation}
均值为$\mu$，方差为$\sigma^2$。标准正态分布的均值为零、方差为一。

多元正态分布$\bm X\sim\mathcal N(\bm\mu,\Sigma)$在$\Sigma\succ0$时具有密度
\begin{equation}
 p(\bm x)=\frac1{(2\pi)^{d/2}|\Sigma|^{1/2}}
 \exp\left[-\frac12(\bm x-\bm\mu)^{\mathsf T}\Sigma^{-1}(\bm x-\bm\mu)\right].
 \label{eq:math-gaussian-density}
\end{equation}
$\bm x,\bm\mu\in\R^d$，$\Sigma\in\R^{d\times d}$为协方差矩阵，$|\Sigma|=\det\Sigma>0$。若$\bm Z$的分量为相互独立的标准正态变量，则$\bm\mu+\Sigma^{1/2}\bm Z$具有该分布。固定$B\in\R^{m\times d}$、$\bm b\in\R^m$时，$B\bm X+\bm b\sim\mathcal N(B\bm\mu+\bm b,B\Sigma B^{\mathsf T})$。

允许$\Sigma\succeq0$时仍可通过线性变换定义退化正态分布，但奇异$\Sigma$没有$\R^d$上的普通密度，不能使用式\eqref{eq:math-gaussian-density}。联合正态变量的零协方差等价于独立，该等价性不适用于任意分布。

设分块联合正态向量为
\begin{equation}
 \begin{pmatrix}\bm X\\\bm Y\end{pmatrix}\sim\mathcal N\left(
 \begin{pmatrix}\bm\mu_X\\\bm\mu_Y\end{pmatrix},
 \begin{pmatrix}\Sigma_{XX}&\Sigma_{XY}\\\Sigma_{YX}&\Sigma_{YY}\end{pmatrix}\right).
\end{equation}
$\bm X\in\R^d$、$\bm Y\in\R^k$，交叉协方差$\Sigma_{XY}\in\R^{d\times k}$且$\Sigma_{YX}=\Sigma_{XY}^{\mathsf T}$。当$\Sigma_{XX}\succ0$时，
\begin{align}
 \bm Y\mid\bm X=\bm x&\sim\mathcal N(\bm\mu_{Y\mid x},\Sigma_{Y\mid X}),\\
 \bm\mu_{Y\mid x}&=\bm\mu_Y+\Sigma_{YX}\Sigma_{XX}^{-1}(\bm x-\bm\mu_X),\\
 \Sigma_{Y\mid X}&=\Sigma_{YY}-\Sigma_{YX}\Sigma_{XX}^{-1}\Sigma_{XY}.
\end{align}
条件协方差是舒尔补，不依赖$\bm x$。令$\bm R=\bm Y-\bm\mu_Y-\Sigma_{YX}\Sigma_{XX}^{-1}(\bm X-\bm\mu_X)$，计算可得$\operatorname{Cov}(\bm R,\bm X)=0$。联合正态性使$\bm R$与$\bm X$独立，从而得到上述条件均值与协方差。

\section{全概率与贝叶斯公式}
\subsection{事件形式}
若$B_i$为两两不交、并集为$\Omega$的有限或可数事件，且各$P(B_i)>0$，全概率公式（law of total probability）为
\begin{equation}
 P(A)=\sum_iP(A\mid B_i)P(B_i).
\end{equation}
它由$A=\bigcup_i(A\cap B_i)$及可数可加性推出。零概率的$B_i$可直接按$P(A\cap B_i)=0$处理，避免定义零概率事件上的条件概率。当$P(A)>0$时，贝叶斯公式（Bayes' theorem）为
\begin{equation}
 P(B_j\mid A)=\frac{P(A\mid B_j)P(B_j)}{\sum_iP(A\mid B_i)P(B_i)}.
\end{equation}
分子是交事件的概率，分母是$P(A)$，保证结果归一化。

\subsection{密度形式}
设$\Theta$与$X$具有联合密度，则
\begin{equation}
 p(x)=\int p(x\mid\theta)p(\theta)\,d\theta,\qquad
 p(\theta\mid x)=\frac{p(x\mid\theta)p(\theta)}{p(x)}.
\end{equation}
离散$\Theta$时使用求和。第二式要求$0<p(x)<\infty$，以几乎处处意义理解。$p(\theta)$称为先验分布（prior distribution），$p(\theta\mid x)$称为后验分布（posterior distribution）。固定$x$后，将$p(x\mid\theta)$视为$\theta$的函数称为似然函数（likelihood function），它不要求在$\theta$上积分为一。$p(x)$称为边缘似然（marginal likelihood），也称证据（evidence）。先验归一化且分母为正、有限时，公式给出归一化的后验。

\section{信息量与分布差异}
比较概率模型时，需要量化单个分布的不确定性，以及两个分布的差异。前者可以用平均信息量描述，后者既可以比较同一结果的概率，也可以考虑不同结果之间的距离。这几类量的定义和性质不同，下面分别说明。

\subsection{熵、交叉熵与相对熵}
离散分布的熵（entropy）为$H(p)=-\sum_xp(x)\log p(x)$，交叉熵（cross entropy）为$H(p,q)=-\sum_xp(x)\log q(x)$。约定$0\log0=0$。若$p(x)>0$且$q(x)=0$，交叉熵为$+\infty$。自然对数对应的信息单位为奈特（nat）。

库尔贝克--莱布勒散度（Kullback--Leibler divergence，KL 散度），又称相对熵，为
\begin{equation}
 \KL(p\|q)=\sum_xp(x)\log\frac{p(x)}{q(x)}.
 \label{eq:math-kl-distance}
\end{equation}
采用同样的零值约定。共同基准测度下的密度使用积分。若$p$不相对于$q$绝对连续，散度为$+\infty$。散度非负，为零当且仅当两分布相同。离散情形由$-\log u\geq1-u$得到
\begin{equation}
 \KL(p\|q)\geq\sum_{p(x)>0}(p(x)-q(x))\geq0.
\end{equation}
连续情形有相应积分证明。KL 散度一般不对称，不满足度量的全部条件。熵与交叉熵均有限时，展开对数比得到
\begin{equation}
 H(p,q)=H(p)+\KL(p\|q),\qquad\KL(p\|q)\geq0.
 \label{eq:math-cross-entropy-kl}
\end{equation}
有限性避免未定义的$\infty-\infty$。连续分布使用微分熵（differential entropy），定义为
\begin{equation}
 h(p)=-\int p(x)\log p(x)\,dx.
\end{equation}
这里$p$为相对于所用坐标下勒贝格测度的密度，要求积分有定义。微分熵可为负，并随坐标尺度改变。

多元正态分布在$\Sigma\succ0$时的微分熵为$\tfrac12\log((2\pi e)^d|\Sigma|)$。在具有相同均值与正定协方差、相关积分有限的密度中，它具有最大微分熵。令$q$为该正态密度，因$\log q$为常数加二次型，相同矩意味着$\E_p\log q=\E_q\log q$。于是$\KL(p\|q)=h(q)-h(p)\geq0$。

\subsection{Wasserstein 距离}
对有限共同支撑集上的分布$p,q$，令$c(a,b)$为支撑集上的度量，$\Gamma(p,q)$为边缘分别是$p,q$的联合质量函数集合。一阶瓦瑟斯坦距离（Wasserstein distance）定义为
\begin{equation}
 W_1(p,q)=\min_{\gamma\in\Gamma(p,q)}\sum_{a,b}\gamma(a,b)c(a,b).
 \label{eq:math-wasserstein-distance}
\end{equation}
约束为$\gamma(a,b)\geq0$、$\sum_b\gamma(a,b)=p(a)$、$\sum_a\gamma(a,b)=q(b)$。$\gamma$也称为耦合（coupling），规定两个边缘取值之间的共同质量。有限支撑使最小值存在，底层$c$为度量时$W_1$也为度量。一般度量空间上将求和改为对耦合测度积分，将最小值改为下确界，并要求有限一阶矩以保证距离有限。

\section{期望的参数微分}
参数变化时，函数值与分布都可能变化。对$\theta\in\R^k$，仅对函数求导一般不能给出完整的期望导数：
\begin{equation}
 \nabla_\theta\E_{X\sim p_\theta}[f_\theta(X)]
 \quad\text{一般不能直接写为}\quad
 \E_{X\sim p_\theta}[\nabla_\theta f_\theta(X)].
 \label{eq:math-expectation-gradient-warning}
\end{equation}
特殊条件下两者可以相等，该式说明不能无条件交换运算。

设$p_\theta$相对于固定测度$\nu$具有共同的参数无关支撑集，且$f_\theta,p_\theta$关于$\theta$可微。若在参数的邻域内，$f_\theta p_\theta$及其各参数偏导具有可积控制函数，则可由支配收敛定理交换微分与积分：
\begin{align}
 \nabla_\theta\int f_\theta(x)p_\theta(x)\,d\nu(x)
 &=\int[p_\theta(x)\nabla_\theta f_\theta(x)
 +f_\theta(x)\nabla_\theta p_\theta(x)]\,d\nu(x)\\
 &=\E_{p_\theta}[\nabla_\theta f_\theta(X)]
 +\E_{p_\theta}[f_\theta(X)\nabla_\theta\log p_\theta(X)].
 \label{eq:math-score-gradient}
\end{align}
第二步另要求共同支撑上密度为正，且各项可积。向量$\nabla_\theta\log p_\theta(x)\in\R^k$称为得分函数（score function），度量固定观测$x$的对数密度怎样随参数变化。将相同的交换条件用于归一化等式$\int p_\theta\,d\nu=1$，得到
\begin{equation}
 \E_{p_\theta}[\nabla_\theta\log p_\theta(X)]=\bm0.
\end{equation}
这里$\bm0\in\R^k$，说明得分在模型自身分布下的均值为零。支撑集随参数变化时需要处理移动边界，不能直接套用式\eqref{eq:math-score-gradient}。

\section{指数族分布}
\subsection{基准测度、自然参数与归一化}
指数族（exponential family）通过固定统计量与参数的内积表达密度。设$\nu$为固定的$\sigma$有限基准测度（base measure），$h(x)\geq0$为参数无关的基准项，$T(x)\in\R^k$为固定可测函数。$\sigma$有限表示空间可被可数个具有有限测度的集合覆盖。分布写作
\begin{equation}
 p_\theta(x)=h(x)\exp\{\theta^{\mathsf T}T(x)-A(\theta)\},\qquad
 A(\theta)=\log\int h(x)e^{\theta^{\mathsf T}T(x)}\,d\nu(x).
 \label{eq:math-exponential-family}
\end{equation}
$p_\theta$为相对于$\nu$的密度，$\theta\in\R^k$称为自然参数（natural parameter），$T$称为充分统计量（sufficient statistic），$A$称为对数配分函数（log-partition function）。$h$不必积分为一。离散变量可使用计数测度，普通连续密度使用勒贝格测度。也可将$h$吸收到测度$h\,d\nu$中。

自然参数空间为
\begin{equation}
 \mathcal H=\left\{\theta\in\R^k:
 0<\int h(x)e^{\theta^{\mathsf T}T(x)}\,d\nu(x)<\infty\right\}.
\end{equation}
只有$\theta\in\mathcal H$时，密度才可归一化。$A$使$\int p_\theta\,d\nu=1$，同时描述参数改变时归一化量的变化。共同支撑由$h(x)>0$确定，要求不随参数变化。

参数集合取为全部$\mathcal H$时称为满指数族（full exponential family）。自然参数集合为$\R^k$中的开集时称为正则指数族（regular exponential family）。使用受限参数映射$\theta=\theta(\eta)$时，关于$\theta$的导数不能直接当作关于$\eta$的导数。

\subsection{充分性与最小表示}
统计量的充分性是指给定该统计量后，随机对象的条件分布不再依赖参数。对具有共同支配测度的分布族，因子分解判据将充分性写为$p_\theta(x)=g_\theta(S(x))b(x)$。式\eqref{eq:math-exponential-family}具有这一形式，所以$T(X)$对$\theta$充分。

若$X_1,\ldots,X_n$相互独立且同分布，其联合密度为
\begin{equation}
 \prod_{i=1}^np_\theta(x_i)=\left(\prod_{i=1}^nh(x_i)\right)
 \exp\left\{\theta^{\mathsf T}\sum_{i=1}^nT(x_i)-nA(\theta)\right\}.
\end{equation}
因此$\sum_iT(X_i)$是联合分布的充分统计量。充分性不保证维度最小，也不保证不同参数对应不同分布。

若不存在非零$\bm a\in\R^k$和常数$b$使$\bm a^{\mathsf T}T(x)=b$在测度$h\,d\nu$下几乎处处成立，则称为最小表示（minimal representation）。存在这样的仿射关系时，统计量含有冗余分量，自然参数可沿相应方向改变而不改变分布。

\subsection{对数配分函数的导数}
在自然参数空间内部，要求一阶、二阶微分可移入积分，且相关矩有限。一个足够的局部条件是：在参数的某个邻域内，$h(x)e^{\theta^{\mathsf T}T(x)}$及其一阶、二阶偏导具有共同的可积控制函数。令$Z(\theta)=e^{A(\theta)}$，则
\begin{align}
 \frac{\partial A}{\partial\theta_j}
 &=\frac1{Z(\theta)}\int T_j(x)h(x)e^{\theta^{\mathsf T}T(x)}\,d\nu(x)
 =\E_\theta[T_j(X)],\\
 \frac{\partial^2A}{\partial\theta_i\partial\theta_j}
 &=\E_\theta[T_i(X)T_j(X)]-\E_\theta[T_i(X)]\E_\theta[T_j(X)].
\end{align}
因此
\begin{equation}
 \nabla A(\theta)=\E_\theta[T(X)]\in\R^k,\qquad
 \nabla^2A(\theta)=\operatorname{Cov}_\theta(T(X))\in\R^{k\times k}.
 \label{eq:math-log-partition-derivatives}
\end{equation}
这些公式将归一化函数的导数与统计量的矩联系起来。海森矩阵半正定，因为$\bm a^{\mathsf T}\nabla^2A\bm a=\Var_\theta(\bm a^{\mathsf T}T(X))\geq0$。最小表示下，非零线性组合不能几乎必然为常数，所以在满足上述条件的内部点，海森矩阵正定。

自然参数空间是凸集。对$\theta_1,\theta_2\in\mathcal H$与$0<t<1$，赫尔德不等式（Hölder inequality）给出$Z(t\theta_1+(1-t)\theta_2)\leq Z(\theta_1)^tZ(\theta_2)^{1-t}$。取对数可得$A$为凸函数。在最小正则表示下，$A$在自然参数空间内部严格凸。边界处导数或矩可能不存在，不能无条件延伸内部结论。

\subsection{得分与信息矩阵}
在上述正则条件下，得分函数为
\begin{equation}
 \nabla_\theta\log p_\theta(x)=T(x)-\nabla A(\theta),
\end{equation}
其期望为零。费舍尔信息矩阵（Fisher information matrix）定义为得分的二阶矩：
\begin{equation}
 \mathcal I(\theta)=\E_\theta\left[
 (\nabla_\theta\log p_\theta(X))(\nabla_\theta\log p_\theta(X))^{\mathsf T}\right]
 =\nabla^2A(\theta).
\end{equation}
它为$k\times k$半正定矩阵，最小正则表示下正定。它描述分布相对于自然参数的局部变化。若$\theta=\theta(\eta)$可微，则信息矩阵按$\mathcal I_\eta=(D\theta)^{\mathsf T}\mathcal I_\theta D\theta$变换，矩阵元素依赖参数坐标。

